\documentclass[aps,prd,twocolumn,superscriptaddress,nofootinbib]{revtex4-2}

\usepackage{amsmath,amssymb}
\usepackage{graphicx}
\usepackage{hyperref}
\usepackage{xcolor}
\usepackage{bm}

% Convenience macros
\newcommand{\HH}{H_0}
\newcommand{\dL}{d_L}
\newcommand{\Mz}{M_z}
\newcommand{\Mbh}{M_\mathrm{bh}}
\newcommand{\Mstellar}{M_\mathrm{stellar}}
\newcommand{\pGW}{p_\mathrm{GW}}
\newcommand{\pdet}{P_\mathrm{det}}
\newcommand{\pgal}{p_\mathrm{gal}}
\newcommand{\Lcat}{L_\mathrm{cat}}
\newcommand{\Lcomp}{L_\mathrm{comp}}
\newcommand{\SNR}{\varrho}
\newcommand{\todo}[1]{\textcolor{red}{\textbf{[TODO: #1]}}}
\newcommand{\pending}[1]{\textcolor{blue}{\textbf{[RESULT PENDING: #1]}}}

\begin{document}

\title{Constraints on the Hubble constant from EMRI dark sirens with LISA using the massive black hole mass}

\author{Jasper Seehofer}
\affiliation{Institute for Theoretical Physics, Ruprecht-Karls University Heidelberg, Philosophenweg 16, 69120 Heidelberg, Germany}

\author{TBD}
\affiliation{TBD}

\date{\today}

\begin{abstract}
Extreme mass ratio inspirals (EMRIs) detected by the Laser Interferometer
Space Antenna (LISA) are promising dark siren sources for measuring the
Hubble constant~$\HH$.  We present a Bayesian inference pipeline that
cross-matches simulated EMRI detections with the GLADE+ galaxy catalog,
comparing two likelihood channels: a standard three-parameter channel using
the luminosity distance and sky position~$(\dL, \phi, \theta)$, and an
extended four-parameter channel that additionally exploits the redshifted
massive black hole mass~$\Mz$ measured from the gravitational waveform.
LISA measures~$\Mz$ to a fractional precision far below the scatter of
any black-hole--galaxy relation, so the mass channel can at best
re-weight candidate hosts within that scatter.
We use Fisher-matrix-based Cram\'er--Rao bounds for a population of
detected EMRIs and the completeness-corrected galaxy catalog likelihood of
Gray~\textit{et~al.}\ [Phys.\ Rev.\ D \textbf{101}, 122001 (2020)].
The reference EMRI forecasts assume galaxy catalogues that may not hold
the typical host; we quantify what remains.
\pending{Rank-1: every $\HH$ uncertainty conditional on and marginalised
over the EMRI population, the selection function and the
black-hole--galaxy relation, for GLADE+ today and for Stage-IV
photometric and pointed-spectroscopy scenarios}
We discuss the impact of catalog
incompleteness, detection probability estimation, and systematic
effects on the posterior.
\end{abstract}

\maketitle

%% ============================================================
\section{Introduction}
\label{sec:intro}
%% ============================================================

% Introduction — constraints on H0 from EMRI dark sirens with LISA using the MBH mass

The Hubble constant $\HH$ sets the present-day expansion rate of the
Universe and anchors the absolute distance scale in cosmology.  Two
mature, independent determinations of $\HH$ are now in significant
tension.  Observations of the cosmic microwave background (CMB) by the
\textit{Planck} satellite, interpreted within a flat $\Lambda$CDM
cosmology, yield $\HH = 67.4 \pm 0.5\;\mathrm{km\,s^{-1}\,Mpc^{-1}}$~\cite{Planck:2018vyg}.
The local distance ladder, calibrated through Cepheid variables and
Type~Ia supernovae by the SH0ES collaboration, gives $\HH = 73.0 \pm
1.0\;\mathrm{km\,s^{-1}\,Mpc^{-1}}$~\cite{Riess:2021jrx}, a
discrepancy exceeding $5\sigma$.  Whether this tension reflects
unaccounted systematics in one or both measurements, or signals new
physics beyond $\Lambda$CDM, remains an open question.  Resolving it
demands independent probes of $\HH$ that do not share the systematic
uncertainties of either the CMB or the distance ladder.

Gravitational-wave (GW) observations provide such an independent
probe.  As recognized by Schutz~\cite{Schutz:1986gp}, a coalescing
compact binary is a \textit{standard siren}: the GW signal encodes the
luminosity distance $\dL$ directly, without any need for a cosmic
distance ladder.  If the redshift $z$ of the source is also known, the
pair $(\dL, z)$ constrains $\HH$ through the distance--redshift
relation.  The first demonstration came with GW170817, the binary
neutron star merger detected by LIGO and Virgo, whose electromagnetic
counterpart identified the host galaxy NGC~4993 and yielded $\HH =
70.0^{+12.0}_{-8.0}\;\mathrm{km\,s^{-1}\,Mpc^{-1}}$~\cite{Abbott:2017xzu}.
With a growing population of bright sirens accompanied by
electromagnetic counterparts, percent-level $\HH$ measurements become
feasible within a decade of advanced detector
operations~\cite{Chen:2017rfc}.

Most compact binary mergers, however, lack an identifiable
electromagnetic counterpart.  These \textit{dark sirens} require a
statistical approach to obtain redshift information.  In the galaxy
catalog method~\cite{Schutz:1986gp,Gray:2019ksv}, one cross-matches
the GW sky localization and distance estimate with a galaxy catalog,
identifying candidate host galaxies within the localization volume.
The likelihood for $\HH$ is then constructed by marginalizing over all
potential hosts, weighted by their probability of being the true
source.  As the number of dark siren events grows, the contributions
from incorrect host galaxies average out, and the posterior on $\HH$
converges to the true value.  A key subtlety is that galaxy catalogs
are incomplete beyond a redshift that depends on the survey depth; a
completeness correction must be included to avoid biasing the $\HH$
inference~\cite{Gray:2019ksv,Gair:2022zsa}.

The Laser Interferometer Space Antenna (LISA)~\cite{LISA:2024hlh}, an
ESA-led space-based gravitational-wave observatory targeting the
millihertz band, is expected to launch in the mid-2030s.  Among its
richest source classes are extreme mass-ratio inspirals (EMRIs): a
stellar-mass compact object ($\mu \sim 10\,M_\odot$) spiraling into a
massive black hole ($\Mbh \sim 10^5$--$10^7\,M_\odot$) in the center
of a galaxy~\cite{Babak:2017tow,LISA:2024hlh}.  An EMRI signal
persists in the LISA band for $\sim\!10^5$ orbital cycles over
months to years, enabling extraordinarily precise parameter estimation.
Sky localization uncertainties of $\mathcal{O}(1\;\mathrm{deg}^2)$
and fractional luminosity distance errors of a few percent are
achievable for high signal-to-noise ratio (SNR) events~\cite{Babak:2017tow}.
Because EMRIs are not expected to produce detectable electromagnetic
radiation, they are natural dark siren sources.  Their superior sky
localization compared to ground-based dark sirens, however, dramatically
reduces the number of candidate host galaxies per event, making the
statistical galaxy catalog method particularly powerful.

Laghi et al.~\cite{Laghi:2021pqk} performed the first comprehensive
study of $\HH$ inference from EMRI dark sirens with LISA, using
Fisher-matrix errors from the population models of
Ref.~\cite{Babak:2017tow} and a mock catalogue complete above
$M_* \sim 10^{10}\,M_\odot$ to $z \lesssim 0.5$; they found $\HH$ to
1.6--3.6\% (90\% credible) for a four-year mission (2.5\% for model M1)
and 1.1--2.6\% for ten years.

In this work, we extend the EMRI dark siren framework by exploiting an
additional observable that LISA measures with remarkable precision: the
redshifted mass of the central massive black hole, $\Mz = \Mbh(1+z)$.
LISA determines $\Mz$ from the EMRI waveform phase evolution to a
fractional accuracy of $\sigma(\Mz)/\Mz \sim 10^{-5}$--$10^{-4}$,
far exceeding the precision of any electromagnetic mass
proxy~\cite{Babak:2017tow}.  Galaxy catalogs, meanwhile, provide
stellar masses $\Mstellar$ for their entries, which can be mapped to
black hole masses through empirical scaling
relations~\cite{Reines:2015moa}.  Weighting EMRI host candidates by
such a relation was proposed as future work by Laghi
et al.~\cite[Sec.~6.4]{Laghi:2021pqk}, and Zhu
et al.~\cite{Zhu:2024qpp} wrote the redshifted-mass comparison into an
EMRI likelihood, using it at least to estimate catalogue completeness;
neither quantifies the gain on $\HH$ with a scattered relation and an
incomplete catalogue.

We present a complete EMRI dark siren inference pipeline that
incorporates Fisher-matrix-based Cram\'er--Rao bounds for all
14~EMRI parameters, the GLADE+ galaxy catalog~\cite{Dalya:2021ewn},
a completeness correction following Gray et
al.~\cite{Gray:2019ksv}, and a detection probability model
accounting for EMRI population and selection effects.
We compare two likelihood
configurations: a baseline ``without MBH mass'' channel that uses only
the GW-measured luminosity distance and sky position $(\dL, \varphi,
\theta)$ to identify candidate hosts, and an extended ``with MBH
mass'' channel that additionally incorporates the redshifted mass
$\Mz$.  The reference forecasts assume catalogues that may not hold the
typical host; we quantify what remains.
\pending{Rank-1: what the mass channel adds to the $\HH$ posterior,
conditional and marginalised}.

The remainder of this paper is organized as follows.
Section~\ref{sec:method} describes the signal model, parameter
estimation framework, dark siren likelihood construction, galaxy
catalog treatment, and detection probability model.
Section~\ref{sec:results} presents our main results, including
detection statistics, the $\HH$ posterior comparison between the two
likelihood channels, and an assessment of systematic effects.
Section~\ref{sec:discussion} discusses the implications and
limitations of our findings.  We conclude in
Sec.~\ref{sec:conclusions}.  Appendix~\ref{app:parameters} tabulates
the EMRI parameter space used in the simulations.
Throughout, we adopt geometrized units with $G = c = 1$ unless
otherwise noted, and quote $\HH$ in units of
$\mathrm{km\,s^{-1}\,Mpc^{-1}}$.


%% ============================================================
\section{Method}
\label{sec:method}
%% ============================================================

\input{sections/method}

%% ============================================================
\section{Results}
\label{sec:results}
%% ============================================================

% Results section — PRD two-column
% REVISED: production run values pending; thesis baselines used as placeholders

Having established the pipeline in Sec.~\ref{sec:method}, we now present
the results of the EMRI dark siren analysis. We begin with the properties
of the simulated detection catalog (Sec.~\ref{sec:detections}), validate
the parameter estimation against known benchmarks
(Sec.~\ref{sec:pe_validation}), and then present the central result:
the $\HH$ posterior with and without the MBH mass channel
(Sec.~\ref{sec:h0_posterior}). We conclude with the impact of
detection probability and catalog completeness corrections
(Sec.~\ref{sec:pdet_completeness}) and a discussion of systematic
effects (Sec.~\ref{sec:systematics}).

%% ============================================================
\subsection{Simulated EMRI detections}
\label{sec:detections}
%% ============================================================

We simulate EMRI events across a grid of seven trial Hubble constant
values $h \in \{0.60, 0.65, 0.70, 0.73, 0.75, 0.80, 0.90\}$, drawing
source parameters from the astrophysical distributions described in
Sec.~\ref{sec:method}. A total of $N_\mathrm{inj} = 165{,}000$ events
are injected, and the SNR is computed for each using the full LISA TDI
response~\cite{Babak:2023lro}. An event is considered detected when its
network $\SNR$ exceeds the threshold $\SNR_\mathrm{thr} = 15$.

Of the $165{,}000$ injected events, $N_\mathrm{det} = 663$ pass the
detection threshold across all $h$-values, corresponding to an overall
detection fraction $f_\mathrm{det} = 4.0 \times 10^{-3}$. The per-$h$
detection yield increases monotonically from
$f_\mathrm{det}(h=0.60) = 2.2 \times 10^{-3}$ to
$f_\mathrm{det}(h=0.90) = 8.1 \times 10^{-3}$, a factor of $3.6$
increase. This trend is expected: larger $\HH$ implies smaller luminosity
distances at fixed redshift, hence higher SNR. Minor non-monotonicity at
$h = 0.70$ is consistent with Poisson noise at the $1\sigma$ level.

All detected events lie at redshift $z < 0.21$, well below the injection
cutoff at $z_\mathrm{cut} = 0.5$, which provides a safety margin of at
least $2.5\times$. The Farr criterion~\cite{Farr:2019rap}
$N_\mathrm{inj}/N_\mathrm{det} \geq 4$ is satisfied for all $h$-values,
with the minimum ratio of $124$ occurring at $h = 0.90$.

For the posterior evaluation, we further require that each detection has a
relative luminosity distance uncertainty
$\sigma_{\dL}/\dL < 0.10$, as estimated from the Cram\'{e}r--Rao
bound. This cut reduces the sample to
\pending{N\_det posterior after $\sigma_{d_L}/d_L < 0.10$ cut from
production run} detections.

\begin{figure}[t]
  \centering
  % \includegraphics[width=\columnwidth]{figures/snr_distribution.pdf}
  \todo{Figure: Left panel --- histogram of $\SNR$ for detected events
  showing the peak near threshold and power-law tail.
  Right panel --- scatter of $\SNR$ vs $\dL$ colored by redshift,
  showing the detection horizon.}
  \caption{Distribution of detected EMRI events. \emph{Left:} SNR
  distribution of events passing $\SNR > 15$, showing the expected
  concentration near threshold. \emph{Right:} SNR versus luminosity
  distance, colored by source redshift. The detection horizon at
  $\dL \approx 1\,\mathrm{Gpc}$ is clearly visible.}
  \label{fig:snr_distribution}
\end{figure}

%% ============================================================
\subsection{Parameter estimation validation}
\label{sec:pe_validation}
%% ============================================================

The Fisher information matrix for each detected EMRI is computed using
forward-difference derivatives of the full LISA TDI waveform
\cite{Vallisneri:2007ev}, yielding Cram\'{e}r--Rao bounds on the
14-dimensional EMRI parameter space. We validate the Fisher matrix
computation against Table~III of Babak
\emph{et~al.}~\cite{Babak:2017tow}, which provides benchmark
Cram\'{e}r--Rao bounds for a fiducial EMRI source. Our results agree with
the benchmark values to within the expected numerical precision of the
finite-difference stencil.

For the dark siren inference, three parameter estimation outputs are
critical: the luminosity distance uncertainty $\sigma_{\dL}$, the sky
localization area $\Delta\Omega$, and the redshifted MBH mass uncertainty
$\sigma_{\Mz}$. The sky localization is computed as
\begin{equation}
  \Delta\Omega = 2\pi |\sin\theta_S| \,
  \sqrt{\sigma_{\theta_S}^2 \, \sigma_{\phi_S}^2
    - \mathrm{cov}(\theta_S, \phi_S)^2} \,,
  \label{eq:sky_localization}
\end{equation}
where $\sigma_{\theta_S}$ and $\sigma_{\phi_S}$ are the
Cram\'{e}r--Rao bounds on the ecliptic co-latitude and longitude,
respectively.

The redshifted MBH mass $\Mz$ is measured with extraordinary precision
by the gravitational waveform: the fractional uncertainty is
$\sigma_{\Mz}/\Mz \lesssim 10^{-6}$ for sources with
$\SNR \geq 15$~\cite{Babak:2017tow}. In contrast, the BH mass of
candidate host galaxies, inferred from the stellar mass via the
$\Mbh$--$\Mstellar$ scaling relation~\cite{Reines:2015moa}, carries a
fractional uncertainty of $\sigma_{\Mbh}/\Mbh \sim 0.1$. This
difference of roughly five orders of magnitude justifies treating
$\Mz$ as exactly known in the likelihood evaluation: the GW mass
measurement acts as an effective delta function when compared against the
broad galaxy mass distribution, i.e.,
\begin{equation}
  p(\Mz^\mathrm{obs} \mid \Mz^\mathrm{true})
  \approx \delta(\Mz^\mathrm{obs} - \Mz^\mathrm{true}) \,.
  \label{eq:delta_mass}
\end{equation}
This approximation eliminates $\Mz$ as a free parameter in the GW
likelihood, so that the mass information enters solely through the
requirement that the candidate host galaxy's BH mass satisfies
$M_\mathrm{bh} = \Mz / (1+z)$.

%% ============================================================
\subsection{$\HH$ posterior: with vs.\ without MBH mass}
\label{sec:h0_posterior}
%% ============================================================

We evaluate the $\HH$ posterior on a grid of $h$-values spanning
$[0.60, 0.90]$ using the dark siren likelihood described in
Sec.~\ref{sec:method}. Figure~\ref{fig:posterior_comparison} shows the
central result of this work: the comparison of the $\HH$ posterior
obtained \emph{without} the MBH mass channel (using $\dL$, $\phi_S$,
$\theta_S$ only) and \emph{with} the MBH mass channel (additionally
using $\Mz$).

\paragraph{Without MBH mass.}
\pending{Rank-1 / production value for $h^\mathrm{(no\,mass)}$}. The
thesis-era baseline ($\pdet = 1$) is withdrawn: it predates the current
pipeline and no surviving result file reproduces it.

\paragraph{With MBH mass.}
\pending{Rank-1 / production value for $h^\mathrm{(mass)}$}. The
thesis-era baseline ($\pdet = 1$) is withdrawn: it predates the current
pipeline and no surviving result file reproduces it.

\paragraph{Mass-channel ablation.}
\pending{Rank-1: $\sigma_h$ with and without the mass term, conditional
and marginalised, per scenario.}

\begin{figure}[t]
  \centering
  % \includegraphics[width=\columnwidth]{figures/h0_posterior_comparison.pdf}
  \todo{Figure: $\HH$ posterior comparison. Blue solid --- without MBH
  mass. Red dashed --- with MBH mass. Vertical gray line at $h = 0.73$
  (injected value). Shaded $1\sigma$ bands.}
  \caption{Posterior on the Hubble constant $\HH$ from
  \pending{N\_det} EMRI dark siren events. The blue curve uses only
  the luminosity distance and sky position; the red curve additionally
  incorporates the redshifted MBH mass. The vertical
  dashed line marks the injected value $h = 0.73$.}
  \label{fig:posterior_comparison}
\end{figure}

\pending{Rank-1: single-event likelihoods $p(d_\mathrm{GW}^i \mid h)$
with and without the mass term (Fig.~\ref{fig:single_event}).}

\begin{figure}[t]
  \centering
  % \includegraphics[width=\columnwidth]{figures/single_event_likelihoods.pdf}
  \todo{Figure: Single-event likelihoods for 4 representative
  detections. Left column: without MBH mass (broad/flat). Right column:
  with MBH mass.}
  \caption{Single-event likelihoods $p(d_\mathrm{GW}^i \mid h)$ for
  four representative EMRI detections at different redshifts.
  \emph{Left column:} using only $\dL$ and sky position --- the
  likelihoods are broad due to the large number of candidate hosts.
  \emph{Right column:} additionally using $\Mz$.}
  \label{fig:single_event}
\end{figure}

\paragraph{Convergence with event number.}
To assess the scaling of the $\HH$ constraint with the number of
detections, we evaluate the posterior on random subsets of
\pending{subset size} events drawn from the full catalog. The
posterior width narrows approximately as $N_\mathrm{det}^{-1/2}$
for both channels, as expected for independent measurements.
\pending{Rank-1: posterior width against the number of events, per
likelihood channel.}

\begin{figure}[t]
  \centering
  % \includegraphics[width=\columnwidth]{figures/posterior_convergence.pdf}
  \todo{Figure: $1\sigma$ width of $\HH$ posterior vs number of events
  for both channels. Dashed line: $N^{-1/2}$ scaling.}
  \caption{Convergence of the $\HH$ posterior width with the number of
  EMRI detections. Points show the median $1\sigma$ width from
  \pending{N\_subsets} random subsets; error bars indicate the
  16th--84th percentile range. The dashed line shows the expected
  $N^{-1/2}$ scaling.}
  \label{fig:posterior_convergence}
\end{figure}

%% ============================================================
\subsection{Detection probability and completeness effects}
\label{sec:pdet_completeness}
%% ============================================================

The selection-corrected results with a realistic $\pdet$ are
\pending{Rank-1 budget}. We now describe the treatment of $\pdet$ and of
the catalog completeness correction.

\paragraph{Detection probability.}
The detection probability $\pdet(\dL, \Mbh, h)$ is estimated via
importance-sampling--weighted kernel density estimation on the injection
campaign (Sec.~\ref{sec:method}), validated with Wilson confidence
interval overlap across $916$ grid bins with zero
Benjamini--Hochberg--adjusted false
discoveries~\cite{Farr:2019rap,Tiwari:2017ndi}.
\pending{Present results with realistic $\pdet$ from production run.}

\paragraph{Catalog completeness.}
The GLADE+ galaxy catalog~\cite{Dalya:2021ewn} is incomplete at redshifts
$z \gtrsim 0.03$. To avoid the systematic bias toward low $\HH$ that
arises when the catalog is denser at lower redshifts (where smaller $h$
maps the observed $\dL$ to a nearer, better-cataloged volume), we
implement the completeness correction of Gray
\emph{et~al.}~\cite{Gray:2019ksv}. The corrected per-event likelihood is
\begin{equation}
  p_i(x_\mathrm{GW} \mid h) =
    f_i \cdot \Lcat^{(i)}
    + (1 - f_i) \cdot \Lcomp^{(i)} \,,
  \label{eq:completeness_correction}
\end{equation}
where $f_i = f(z_i, h) \in [0, 1]$ is the completeness fraction of the
GLADE+ catalog at the inferred redshift $z_i = z(\dL^{(i)}, h)$, and
$\Lcat^{(i)}$ and $\Lcomp^{(i)}$ are the catalog and completion
likelihood terms, respectively. The completion term accounts for the
probability of a host galaxy not present in the catalog:
\begin{equation}
  \Lcomp^{(i)} =
    \frac{\displaystyle\int dz \; \pGW(x \mid z, \Omega_\mathrm{det}, h)
      \; \pdet(\dL(z,h)) \; \frac{dV_c}{dz}}
    {\displaystyle\int dz \; \pdet(\dL(z,h))
      \; \frac{dV_c}{dz}} \,.
  \label{eq:completion_term}
\end{equation}

The completeness fraction $f(z, h)$ is interpolated from the GLADE+
completeness data. At $h = 0.73$, the catalog is approximately $50\%$
complete at $z \approx 0.03$, dropping to $\sim 35\%$ at
$z \approx 0.11$ and $\sim 21\%$ at $z \gtrsim 0.18$, where it
plateaus due to extrapolation beyond the available data. The $h$-dependence
arises because larger $\HH$ maps a given $z$ to a smaller $\dL$, at which
the catalog is more complete: $f(z = 0.05,\, h = 0.86) > f(z = 0.05,\, h = 0.60)$.

The combination formula~\eqref{eq:completeness_correction} satisfies two
important limiting cases: $f_i = 1$ exactly recovers the catalog-only
likelihood, and $f_i = 0$ gives the pure completion term, which uses only
the comoving volume prior and the GW likelihood. These limits, together
with the requirement that the combined likelihood interpolates
monotonically between the two terms, have been verified analytically and
numerically.

\pending{Show the effect of the completeness correction on the $\HH$
posterior bias (production run, $h_\mathrm{true} = 0.73$).}

%% ============================================================
\subsection{Systematic effects}
\label{sec:systematics}
%% ============================================================

We identify three dominant sources of systematic uncertainty in the
analysis.

\paragraph{Galaxy catalog incompleteness.}
As discussed in Sec.~\ref{sec:pdet_completeness}, the primary source of
systematic bias is the incompleteness of the GLADE+ catalog at
$z \gtrsim 0.03$. Without the completeness correction, the posterior is
systematically pulled toward lower $\HH$ because the catalog is
preferentially populated at low redshifts, where a smaller $h$ can
reproduce the observed $\dL$. The magnitude of this effect is
\pending{Rank-1 budget}. The completeness
correction of Eq.~\eqref{eq:completeness_correction} mitigates this bias,
though its effectiveness depends on the accuracy of the completeness
function $f(z, h)$, which is currently based on digitized data from
Ref.~\cite{Dalya:2021ewn}.

\paragraph{Number of candidate host galaxies.}
Without the MBH mass channel, each GW detection admits $\sim 10$--$10^4$
candidate host galaxies within the $3\sigma$ sky localization and distance
error volume. The true host contribution to the single-event likelihood
can be as small as $\sim 0.6\%$ for high-redshift detections
($z \gtrsim 0.12$), where the galaxy density is high and the distance
measurement is less constraining.
\pending{Rank-1 census of candidates per event with and without the mass
window, GLADE+ today}.

\paragraph{Detection probability normalization.}
The detection probability $\pdet$ enters both the numerator and
denominator of the single-event likelihood. An $h$-dependent bias in the
$\pdet$ estimate can propagate into a systematic shift of the posterior.
In the $\pdet = 1$ baseline, this effect is absent by construction, but
with a realistic $\pdet$ estimated from an injection campaign of
finite size, the normalization carries a statistical uncertainty of
$\sim 1/\sqrt{N_\mathrm{inj}}$ per grid bin.
\pending{Quantify residual $\pdet$ normalization effect from production
run.}

Table~\ref{tab:results_summary} summarizes the key results for both
inference channels.

\begin{table}[t]
  \caption{Summary of $\HH$ constraints from EMRI dark siren inference.
  The injected value is $h_\mathrm{true} = 0.73$.
  \pending{Update with production run values.}}
  \label{tab:results_summary}
  \centering
  \begin{tabular}{l c c}
    \hline\hline
    & Without $\Mz$ & With $\Mz$ \\
    \hline
    MAP estimate &
      \pending{X.XXX} &
      \pending{X.XXX} \\
    $1\sigma$ width &
      \pending{0.0XX} &
      \pending{0.0XX} \\
    Precision (\%) &
      \pending{X.X} &
      \pending{X.X} \\
    Bias $\Delta h / h_\mathrm{true}$ (\%) &
      \pending{X.X} &
      \pending{X.X} \\
    $N_\mathrm{det}$ used &
      \multicolumn{2}{c}{\pending{NNN}} \\
    \hline\hline
  \end{tabular}
\end{table}

%% CITATIONS NEEDED (for gpd-bibliographer):
%% None — all citations used in this section are already in references.bib.


%% ============================================================
\section{Discussion}
\label{sec:discussion}
%% ============================================================

% discussion.tex — Section IV: Discussion
% Target: PRD two-column (revtex4-2)
% Macros used: \HH, \dL, \Mz, \Mbh, \Mstellar, \pGW, \pdet, \Lcat, \Lcomp, \SNR, \pending, \todo

We now place the analysis in the context of prior work, discuss what
the MBH mass can and cannot add, and assess the limitations of the
analysis.

%% ============================================================
\subsection{Comparison with previous work}
\label{sec:comparison}
%% ============================================================

The most direct comparison is with
Laghi~\emph{et~al.}~\cite{Laghi:2021pqk}, who performed the first
comprehensive study of EMRI dark sirens as probes of $\HH$ with LISA.
Our analysis extends their work in four significant ways.

First, Laghi~\emph{et~al.}\ took Fisher-matrix errors from the
population models of Ref.~\cite{Babak:2017tow}; we compute the full
Cram\'{e}r--Rao bound from the $14 \times 14$ Fisher information matrix
for every simulated event (Sec.~\ref{sec:fisher}), which captures the
dependence of $\sigma_{\dL}$ on SNR, sky position, mass ratio, and spin.
The resulting per-event distance uncertainties span a wide range, with
some high-SNR detections achieving $\sigma_{\dL}/\dL < 1\%$.

Second, we introduce the MBH mass channel, which exploits the precise
measurement of $\Mz$ from the gravitational waveform as an additional
discriminant for host galaxy identification.  The idea is not new:
Laghi~\emph{et~al.}\ proposed it as future work
(Ref.~\cite{Laghi:2021pqk}, Sec.~6.4), and
Zhu~\emph{et~al.}~\cite{Zhu:2024qpp} included a term comparing $\Mz$
with $(1+z)\hat{M}(L_K)$ in their EMRI likelihood
($0.4\,\mathrm{dex}$, window $[-3\sigma, +2\sigma]$), used at least for
completeness and without an ablation of the mass term.  What this work
adds is that ablation, mass weighting on and off, on a real catalogue
(GLADE+) \pending{Rank-1: with the relation's scatter drawn in the mock
and its normalisation and scatter marginalised}.

Third, we employ the real GLADE+ galaxy catalog~\cite{Dalya:2021ewn},
which includes spectroscopic and photometric redshifts, stellar masses,
and sky positions for approximately 23~million objects.
Laghi~\emph{et~al.}\ used a synthetic catalog with idealized properties.
The real catalog introduces the practical complication of spatially
varying completeness, which we address by implementing the correction of
Gray~\emph{et~al.}~\cite{Gray:2019ksv}.

Fourth, we include the catalog completeness correction
[Eq.~\eqref{eq:completeness_combination}] following Gray~\emph{et~al.},
which was originally developed for ground-based dark siren analyses of
binary black hole mergers~\cite{Gray:2019ksv}.  This correction is
essential for avoiding the systematic low-$\HH$ bias caused by
preferentially cataloged low-redshift galaxies, as demonstrated in
Sec.~\ref{sec:systematics}.

Our Fisher matrix results are validated against the benchmark
Cram\'{e}r--Rao bounds in Table~III of
Babak~\emph{et~al.}~\cite{Babak:2017tow}, finding agreement to within
the numerical precision of the finite-difference stencil
(Sec.~\ref{sec:pe_validation}).  This provides confidence that the
parameter estimation inputs to the dark siren likelihood are correctly
computed.

The completeness correction we implement follows the framework of
Gray~\emph{et~al.}~\cite{Gray:2019ksv}, adapted here for EMRI sources
with LISA.  Gair~\emph{et~al.}~\cite{Gair:2022zsa} recently provided a
comprehensive study of systematic biases in dark siren analyses,
identifying catalog incompleteness as the dominant source of error for
events beyond $z \sim 0.03$ in GLADE+.  Our treatment of the
completeness correction is consistent with their recommendations,
although we note that our completeness function is interpolated from
digitized data rather than derived from a self-consistent luminosity
function model (see Sec.~\ref{sec:caveats}).

%% ============================================================
\subsection{What the MBH mass can and cannot add}
\label{sec:mass_interpretation}
%% ============================================================

The mass channel's information is limited by the scatter of the
black-hole--galaxy relation, not by the precision of $\Mz$. The
thesis-era factor $\sim 2$ is withdrawn.
\pending{Rank-1: the gain on $\sigma_{\HH}$ at the relation's scatter,
per event and in total}.

Without the mass channel, a GW detection constrains $\dL$ (and sky
position) but not $z$ independently.  For a given trial value of $\HH$,
the inferred redshift $z = z(\dL, \HH)$ identifies a shell in the galaxy
catalog.  All galaxies consistent with the measured distance and sky
position contribute to the single-event likelihood, potentially
numbering $\sim 10$--$10^4$ candidates depending on the sky
localization volume and local galaxy density.  Because each candidate is
weighted equally (up to the GW measurement uncertainty), the per-event
$\HH$ constraint is diluted by the large number of viable hosts.

With the mass channel, the precisely measured $\Mz$ imposes an
additional constraint: the true host must satisfy
$\Mbh \approx \Mz / (1+z)$, where $z$ is determined by $\dL$ and the
trial $\HH$.  Since the GW measurement determines $\Mz$ to a fractional
precision of $\lesssim 10^{-6}$ (Sec.~\ref{sec:pe_validation}), the
width of the mass filter is set by the scatter of the
$\Mbh$--$\Mstellar$ relation (see below), not by the GW measurement.
\pending{Rank-1 census: how many candidates per event the mass window
removes, and how often it removes the true host}.

Any gain from the mass channel should be largest for detections in regions of high galaxy
density (large angular extent and moderate distance), where the standard
channel admits the most candidate hosts.  Conversely, for isolated
detections in underdense regions where only a handful of galaxies lie
within the error volume, the mass information adds less because the host galaxy is already effectively
identified by proximity.

We note that the efficacy of the mass channel is limited by the accuracy
of the $\Mbh$--$\Mstellar$ scaling relation.  The
Reines~\&~Volonteri~\cite{Reines:2015moa} relation has an intrinsic
scatter of approximately $0.55\,\mathrm{dex}$, with best-fit parameters
$\alpha = 7.45 \pm 0.08$ and $\beta = 1.05 \pm 0.11$.  Improvements in
the empirical calibration of this relation, particularly in the
$10^5$--$10^6\,M_\odot$ range relevant for EMRI sources, would directly
translate into a tighter mass filter and thus a stronger $\HH$
constraint.

%% ============================================================
\subsection{Limitations and caveats}
\label{sec:caveats}
%% ============================================================

We identify several approximations and limitations of the present
analysis.  We discuss them in roughly decreasing order of expected
impact on the results.

\paragraph{Fisher matrix approximation.}
The parameter estimation relies on the Fisher information matrix, which
provides a lower bound on parameter variances that is tight only in the
high-SNR, unimodal posterior regime~\cite{Vallisneri:2007ev}.  While our
detection threshold of $\SNR \geq 20$ places the analysis in a regime
where the Fisher approximation is generally reliable, multimodal
posteriors can arise for EMRI sources due to correlations between mass,
spin, and orbital parameters.  In such cases, the Fisher matrix may
underestimate the true parameter uncertainty, leading to overly
optimistic Cram\'{e}r--Rao bounds.  A full Bayesian parameter estimation
using, e.g., Markov chain Monte Carlo (MCMC) sampling would capture
these effects but is computationally prohibitive for the $\sim 10^5$
injections considered here.

\paragraph{Finite-difference accuracy.}
The waveform derivatives entering the Fisher matrix are computed using
forward finite differences with step size $\epsilon = 10^{-6}$,
yielding $\mathcal{O}(\epsilon)$ accuracy.  The five-point stencil of
Eq.~\eqref{eq:five_point}, which achieves $\mathcal{O}(\epsilon^4)$
accuracy, is implemented in the codebase but was not used for the
injection campaign due to the $4\times$ increase in waveform evaluations
per parameter (56 vs.\ 14 per Fisher matrix).  While the forward-difference
and five-point results agree for the validation source of
Ref.~\cite{Babak:2017tow}, the higher-order stencil would provide
improved robustness for sources near the edges of parameter space.

\paragraph{Simplified population model.}
We draw EMRI source parameters from a random distribution over the
parameter ranges listed in Appendix~\ref{app:parameters}, rather than
from an astrophysically motivated population model.  Realistic EMRI rate
estimates~\cite{Babak:2017tow} predict event rates that depend strongly
on the MBH mass function, the stellar environment of the galactic
nucleus, and the efficiency of orbital capture.  A population-informed
injection campaign may yield a different detection yield and a different
distribution of per-event constraints, which could affect both the
overall $\HH$ precision and the relative improvement from the mass
channel.

\paragraph{$\Mbh$--$\Mstellar$ scaling relation.}
The mass channel relies on the empirical scaling relation of
Reines~\&~Volonteri~\cite{Reines:2015moa} to infer the MBH mass from
the stellar mass of candidate host galaxies.  This relation was
calibrated in the local universe ($z \lesssim 0.055$) and may evolve
with redshift.  Moreover, the intrinsic scatter of the relation
($\sim 0.55\,\mathrm{dex}$) is the dominant source of uncertainty in
the galaxy-side mass estimate.  Systematic shifts in the scaling
relation parameters $\alpha$ and $\beta$ would propagate directly into
the mass filter, potentially biasing or broadening the $\HH$ posterior.
Future work should marginalize over the scaling relation parameters
within their uncertainties.

\paragraph{GLADE+ catalog completeness.}
The completeness fraction $f(z, h)$ entering the correction of
Eq.~\eqref{eq:completeness_combination} is interpolated from published
completeness data for the GLADE+ catalog~\cite{Dalya:2021ewn}.  We note
two caveats.  First, the published completeness may reflect
\emph{number} completeness (fraction of galaxies detected) rather than
\emph{luminosity} completeness (fraction of total luminosity captured),
which is the more physically relevant quantity for galaxy-catalog-based
analyses~\cite{Gray:2019ksv}.  Second, the completeness is treated as a
function of redshift alone, whereas the actual catalog completeness
varies with sky position, galaxy type, and apparent magnitude.  A more
refined treatment would use the angular-dependent completeness, though
this is beyond the scope of the present work.

\paragraph{Cosmological model.}
We assume a flat $\Lambda$CDM cosmology with fixed matter density
$\Omega_m = 0.25$.  This value reflects the WMAP-era best
fit~\cite{Hinshaw:2012aka} rather than the Planck~2018 result
$\Omega_m = 0.3153 \pm 0.0073$~\cite{Planck:2018vyg}.  While the
choice of $\Omega_m$ affects the $\dL$--$z$ relation and hence the
inferred $\HH$, the sensitivity of the posterior to $\Omega_m$ is
subdominant to the statistical uncertainty from the finite number of
detections.  A joint inference over $(\HH, \Omega_m)$ or an extension
to the $w$CDM model is straightforward with the existing pipeline
infrastructure but is deferred to future work.

\paragraph{Galactic confusion noise.}
The LISA noise PSD used in this analysis
[Eq.~\eqref{eq:psd}] does not include the foreground from unresolved
galactic white-dwarf binaries, which dominates the sensitivity in the
$0.1$--$3\,\mathrm{mHz}$ band.  While the noise model includes the
$S_c(f)$ term in the formal expression, this component is set to zero in
the current implementation.  Including the confusion noise would reduce
the effective SNR for EMRI sources with significant power in this band,
potentially lowering the detection yield and increasing the per-event
distance uncertainty.  The effect is expected to be moderate for the MBH
mass range considered here ($10^5$--$10^6\,M_\odot$), whose peak GW
emission lies at frequencies above the confusion noise band.

%% ============================================================
\subsection{Future directions}
\label{sec:future}
%% ============================================================

Several extensions of this work are planned or in progress.

\paragraph{Completeness-corrected production run.}
The selection-corrected results are \pending{Rank-1 budget}.  A production-scale simulation incorporating the
completeness-corrected likelihood of
Eq.~\eqref{eq:completeness_combination} with a realistic $\pdet$
estimated from the injection campaign is underway.  This will quantify
the bias reduction from the Gray~\emph{et~al.}\ correction and
establish the definitive $\HH$ constraint.

\paragraph{Astrophysical population model.}
Replacing the uniform parameter draws with an astrophysically motivated
EMRI population model~\cite{Babak:2017tow} --- incorporating the MBH
mass function, stellar cusp properties, and orbital capture
efficiencies --- would provide a more realistic estimate of both the
detection yield over the LISA mission lifetime and the expected $\HH$
precision.

\paragraph{Full Bayesian parameter estimation.}
MCMC-based parameter estimation for individual EMRI detections would
replace the Fisher matrix approximation and correctly capture multimodal
and non-Gaussian posterior features.  While this is computationally
expensive for the full injection campaign, targeted MCMC runs on a
representative subset of detections would validate the Fisher-matrix
results and quantify any systematic bias from the Gaussian
approximation.

\paragraph{Five-point stencil derivatives.}
Upgrading the waveform derivatives from forward differences to the
five-point stencil of Eq.~\eqref{eq:five_point} for the full injection
campaign would improve the Fisher matrix accuracy at the cost of a
$4\times$ increase in waveform evaluations.  This is straightforward to
implement (the five-point stencil is already available in the codebase)
and would be worthwhile for a production-quality analysis.

\paragraph{Joint analysis with other LISA sources.}
EMRIs are one of several source classes observable by LISA.  Combining
the EMRI dark siren constraints with those from massive black hole
binaries (MBBHs) and stellar-origin binary black holes (SOBBHs) would
improve the overall $\HH$ precision, particularly if the different
source classes probe complementary redshift
ranges~\cite{LISA:2024hlh,Chen:2017rfc}.

\paragraph{Dark energy equation of state.}
The pipeline infrastructure already supports a $w$CDM cosmology with
parameters $(w_0, w_a)$, though this extension has not yet been
exercised.  With a sufficiently large detection catalog, EMRI dark
sirens could provide independent constraints on the dark energy equation
of state, complementary to those from Type~Ia
supernovae~\cite{DES:2024tys} and baryon acoustic
oscillations.

\paragraph{Improved completeness modeling.}
An angular-dependent and luminosity-weighted completeness model for
GLADE+ would improve the fidelity of the correction in
Eq.~\eqref{eq:completeness_combination}.  In combination with upcoming
spectroscopic surveys (e.g., DESI, Euclid), which will substantially
increase the catalog depth at $z \gtrsim 0.1$, the completeness
correction is expected to become less critical as the catalog approaches
full coverage over the EMRI detection volume.

%% CITATIONS NEEDED (for gpd-bibliographer):
%% - MISSING:wmap-9yr-2013 — Hinshaw et al. (2013), WMAP 9-year cosmological parameters, arXiv:1212.5226
%% - MISSING:des-y5-2024 — DES Collaboration, Year 5 supernova cosmology results


%% ============================================================
\section{Conclusions}
\label{sec:conclusions}
%% ============================================================

% conclusions.tex — Section V: Conclusions
% Target: PRD two-column (revtex4-2)
% Macros used: \HH, \dL, \Mz, \Mbh, \Mstellar, \pdet, \pending

We have investigated whether the redshifted massive black hole mass
$\Mz$, precisely measured from the EMRI gravitational waveform, can
improve dark siren constraints on the Hubble constant $\HH$ with LISA.
The thesis-era answer, a factor $\sim 2$, is withdrawn: no surviving
result reproduces it.  The reference forecasts assume catalogues that
may not hold the typical host; we quantify what remains
\pending{Rank-1 summary}.

As a methodological contribution, we have adapted the catalog
completeness correction of
Gray~\emph{et~al.}~\cite{Gray:2019ksv} --- originally developed for
ground-based dark siren analyses of binary black hole mergers --- to
the EMRI context with LISA and the GLADE+ galaxy
catalog~\cite{Dalya:2021ewn}.  This correction mitigates the
systematic low-$\HH$ bias caused by the preferential cataloging of
low-redshift galaxies (Sec.~\ref{sec:pdet_completeness}).

Several directions for future work will strengthen these results.
A production-scale simulation incorporating the full detection
probability $\pdet(\dL, \Mbh, h)$ and the completeness-corrected
likelihood is underway and will establish the definitive $\HH$
constraint.  Replacing the uniform EMRI population model with an
astrophysically motivated rate model~\cite{Babak:2017tow} will
provide more realistic detection yields.  Finally, marginalizing over
the parameters of the $\Mbh$--$\Mstellar$ scaling
relation~\cite{Reines:2015moa} and extending the analysis to a joint
$(\HH, \Omega_m)$ or $w$CDM inference are straightforward with the
existing pipeline infrastructure.


%% ============================================================
\begin{acknowledgments}
\todo{Add acknowledgments: advisors, computing resources (bwUniCluster), funding}
\end{acknowledgments}

\bibliography{references}

%% ============================================================
\appendix

\section{EMRI parameter space}
\label{app:parameters}

% appendix_parameters.tex — Appendix A: EMRI Parameter Space
% Target: PRD two-column (revtex4-2)
% Macros used: \HH, \dL, \Mz, \Mbh, \Mstellar

Table~\ref{tab:emri_parameters} lists the 15 parameters characterizing
each simulated EMRI event.  Source parameters that depend on the
astrophysical context --- sky position, luminosity distance, and MBH
mass --- are drawn from the GLADE+ galaxy catalog~\cite{Dalya:2021ewn}.
The MBH mass is inferred from the catalog stellar mass via the
$\Mbh$--$\Mstellar$ scaling relation of
Reines~\&~Volonteri~\cite{Reines:2015moa}, restricted to the range
$10^5$--$10^6\,M_\odot$ consistent with EMRI formation channels.
Extrinsic parameters (sky angles and orbital phases) are drawn uniformly
over their natural ranges.  The compact object mass $\mu = 10\,M_\odot$
and spin $a = 0.98$ are held fixed, following the fiducial M1 model of
Babak~\emph{et~al.}~\cite{Babak:2017tow}.  The initial semi-latus
rectum $p_0$ is restricted from above to ensure the orbit lies outside
the separatrix of the augmented analytic kludge (AAK)
waveform model~\cite{Chua:2020stf,Katz:2021yft}.

We emphasize that the population model adopted here --- uniform draws
in mass, eccentricity, and inclination --- is simplified compared
to astrophysical EMRI rate models~\cite{Babak:2017tow}.  This choice
is deliberate: the goal of this work is to quantify how much of the $\HH$
information survives once the population, the selection and the
black-hole--galaxy relation are marginalised, for which a simplified but controlled population is
sufficient.  The effect of a realistic population model on the absolute
$\HH$ precision is discussed in Sec.~\ref{sec:caveats}.

\begin{table}[t]
  \caption{EMRI parameter space used in the simulation campaign.
  ``Catalog'' indicates the value is drawn from the GLADE+ galaxy
  catalog~\cite{Dalya:2021ewn}; ``Uniform'' denotes a uniform draw
  over the stated range; ``Fixed'' parameters are held constant.
  The redshifted MBH mass $\Mz$ is derived from $\Mbh$ and the
  catalog redshift $z$.}
  \label{tab:emri_parameters}
  \centering
  \begin{tabular}{l l l l}
    \hline\hline
    Parameter & Symbol & Range / Value & Source \\
    \hline
    MBH mass & $\Mbh$ &
      $[10^5, 10^6]\,M_\odot$ & Catalog \\
    Compact object mass & $\mu$ &
      $10\,M_\odot$ & Fixed \\
    Spin parameter & $a$ &
      $0.98$ & Fixed \\
    Initial eccentricity & $e_0$ &
      $[0, 0.2]$ & Uniform \\
    Semi-latus rectum & $p_0$ &
      $[10, 16]\,M$ & Uniform \\
    Inclination cosine & $x_{I,0} = \cos I$ &
      $[-1, 1]$ & Uniform \\
    Luminosity distance & $\dL$ &
      $\dL(z, h_\mathrm{true})$ & Catalog \\
    Sky colatitude & $\theta_K$ &
      $[0, \pi]$ & Uniform \\
    Sky longitude & $\phi_K$ &
      $[0, 2\pi]$ & Uniform \\
    Initial orbital phase & $\Phi_{\varphi,0}$ &
      $[0, 2\pi]$ & Uniform \\
    Initial polar phase & $\Phi_{\theta,0}$ &
      $[0, 2\pi]$ & Uniform \\
    Initial radial phase & $\Phi_{r,0}$ &
      $[0, 2\pi]$ & Uniform \\
    Observation time & $T_\mathrm{obs}$ &
      $5\,\mathrm{yr}$ & Fixed \\
    Sampling interval & $\Delta t$ &
      $10\,\mathrm{s}$ & Fixed \\
    Redshifted MBH mass & $\Mz = \Mbh(1+z)$ &
      Derived & --- \\
    \hline\hline
  \end{tabular}
\end{table}


\end{document}
